% Research component for root integration. All labels have the dq prefix.
% Bibliography required: OlverShakiban1978; incoming canonical fifth report.
% This is a quotient of VECTOR SPACES, not a quotient algebra by an ideal.

\section{Canonical integration by parts for polygamma polynomials}
\label{dq:section}

The incoming canonical fifth-moment report asks for an all-degree
classification of differential monomials after integration by parts,
with their finite endpoint constants retained. The next theorem answers
the formal integration-by-parts part of that question. It also gives an
exact reduction algorithm and a smaller parity sector for the inverse-Gamma
polynomials. The vector-space dimensions below do not count arithmetically
independent constants, or independent images under any particular spectral
map. Additional special-function identities can reduce the resulting
integrals further; the sixth-degree calculation below does exactly this.

Differential-polynomial methods and the exactness of the Euler-operator
complex are established tools; see Olver--Shakiban
\cite{OlverShakiban1978}. The elementary construction below is included
in full to fix the particular grading, basis, and endpoint convention
needed for this application.

\subsection{The formal vector-space quotient and a triangular basis}

Let
\[
 \mathscr R=\mathbb Q[X_1,X_2,\ldots],\qquad
 D=\sum_{j\ge1}X_{j+1}\frac{\partial}{\partial X_j}.
\]
Each polynomial involves finitely many variables, so this derivation
is well defined. Give $X_j$ weight $j$ and factor degree one. Write
$\mathscr R_{N,\ell}$ for the span of monomials of weight $N$ and
factor degree $\ell$, and $\mathscr R_N=\bigoplus_\ell\mathscr R_{N,\ell}$.
The derivation raises weight by one and preserves factor degree.
The substitution $X_j=\psi^{(j-1)}(x)$ intertwines $D$ with $d/dx$.

For a partition $\lambda=(\lambda_1\ge\cdots\ge\lambda_\ell>0)$,
put $X_\lambda=\prod_{i=1}^\ell X_{\lambda_i}$.
Let $\mathscr B_N$ consist of the partitions of $N$ whose largest
part occurs at least twice. At weights zero and one, use separately
the representatives $1$ and $X_1$.

\begin{theorem}[Canonical differential-polynomial normal form]
\label{dq:normalform}
For $N\ge2$, the classes
\[
 \{X_\lambda:\lambda\in\mathscr B_N\}
\]
are a basis of the vector space $\mathscr R_N/D\mathscr R_{N-1}$.
Every $F\in\mathscr R_N$ has a unique decomposition
\begin{equation}\label{dq:decomposition}
 F=\mathcal N(F)+D\mathcal T(F),\qquad
 \mathcal N(F)\in\operatorname{span}_{\mathbb Q}
                 \{X_\lambda:\lambda\in\mathscr B_N\},
 \quad\mathcal T(F)\in\mathscr R_{N-1}.
\end{equation}
If $p_\ell(N)$ is the number of partitions of $N$ into exactly
$\ell$ parts, then
\begin{equation}\label{dq:dimensions}
 \dim \frac{\mathscr R_{N,\ell}}{D\mathscr R_{N-1,\ell}}
   =p_\ell(N)-p_\ell(N-1),\qquad
 \dim \frac{\mathscr R_N}{D\mathscr R_{N-1}}
   =p(N)-p(N-1).
\end{equation}
The decomposition is produced by a terminating rational algorithm.
Every retained monomial uses derivatives of order at most
$\lfloor N/2\rfloor-1$ after substitution into $\psi$.
\end{theorem}

\begin{proof}
Order partitions lexicographically by their decreasing lists of parts.
Let a monomial have a unique largest part $M\ge2$. Decrease this
part to $M-1$, and call the resulting partition $\mu$. If $b$ is
the multiplicity of $M-1$ in $\mu$, the product rule gives
\begin{equation}\label{dq:pivot}
 D X_\mu=bX_\lambda+R_\lambda.
\end{equation}
Every monomial in $R_\lambda$ has largest part at most $M-1$:
it arises by differentiating one of the smaller factors of $X_\mu$.
Thus every such partition is lexicographically smaller than $\lambda$.
Replace $X_\lambda$ by $b^{-1}DX_\mu-b^{-1}R_\lambda$.
Applying this replacement to the largest remaining unique-largest
partition strictly decreases a finite ordered set of partitions of $N$,
so it terminates in the asserted remainder. The replacement preserves
factor degree.

For independence, apply $D$ to the partition monomials of weight $N-1$.
The leading partition is obtained by increasing one largest part by one;
its coefficient is the multiplicity of that largest part and is nonzero.
This is a bijection between partitions of $N-1$ and partitions of $N$
with a unique largest part. It also preserves the number of parts.
Distinct derivative columns therefore have distinct leading partitions.
No nonzero linear combination of those columns can be supported entirely
on repeated-largest partitions. This proves both independence and
injectivity of $D$ on $\mathscr R_{N-1}$, and hence uniqueness of
the primitive in \eqref{dq:decomposition}. Counting the columns and
rows gives \eqref{dq:dimensions}. Finally $2\lambda_1\le N$ for
each retained partition, giving the derivative-order bound.
\end{proof}

The quotient is a vector-space quotient: $D\mathscr R$ is not in
general an ideal, so multiplication of classes is not being asserted.
The ordinary generating series, including weights zero and one, is
\begin{equation}\label{dq:Hilbert}
 \sum_{N\ge0}\dim(\mathscr R_N/D\mathscr R_{N-1})q^N
 =q+\prod_{j\ge2}\frac1{1-q^j}.
\end{equation}
The convention is $D\mathscr R_{-1}=0$. The extra $q$ records
$X_1$, since differentiation kills the constant polynomial.
With factor degree recorded by $z$, the corresponding series is
\[
 (1-q)\prod_{j\ge1}(1-zq^j)^{-1}+q.
\]

For example, the respective bases at weights four, five, and six are
\[
 \begin{array}{c|l}
 4& X_1^4,\ X_2^2\\
 5& X_1^5,\ X_1X_2^2\\
 6& X_1^6,\ X_1^2X_2^2,\ X_2^3,\ X_3^2.
 \end{array}
\]
This explains formally why the fifth digamma calculation introduces
$\psi(\psi')^2$ after the scalar power. It does not assert that its
integral is independent of the spectral constants already present.

\subsection{The exact endpoint functional}

For a meromorphic function with a finite principal part at zero, use
the original-coordinate finite part
\[
 \operatorname{FP}\int_0^1f(x)\,dx
 =\operatorname{CT}_{\varepsilon\downarrow0}\int_\varepsilon^1 f(x)\,dx.
\]
Here the constant term is taken after all negative powers of
$\varepsilon$ and all nonconstant powers of $\log\varepsilon$ are removed.
For every differential polynomial $T$,
\begin{equation}\label{dq:endpointfunctional}
 \operatorname{FP}\int_0^1 DT(\psi,\psi',\ldots)\,dx
 =\mathcal E(T):=T(\psi(1),\psi'(1),\ldots)
          -\operatorname{CT}_{x=0}T(\psi(x),\psi'(x),\ldots).
\end{equation}
This follows by the fundamental theorem of calculus before the constant
term is taken. In particular a total derivative cannot simply be
discarded in the finite-part problem.

The functional is an explicitly finite polynomial calculation. At
weight $w$, use only
\begin{equation}\label{dq:endpointtruncation}
 U_w(x)=-x^{-1}-\gamma+
       \sum_{k=1}^{w-1}(-1)^{k+1}\zeta(k+1)x^k,
\end{equation}
and substitute $X_j=U_w^{(j-1)}(x)$ to obtain the constant term.
Terms omitted from $\psi$ start at $x^w$; in a weight-$w$ monomial,
any occurrence of an omitted term has positive total exponent even
when all other factors contribute their leading poles. Thus it cannot
affect the constant coefficient. The endpoint at one uses
\[
 X_1=-\gamma,\qquad X_j=(-1)^j(j-1)!\zeta(j)\quad(j\ge2).
\]
Consequently $\mathcal E(T)\in\mathbb Q[\gamma,\zeta(2),\ldots,\zeta(w)]$
for every $T\in\mathscr R_w$.

For a formula with entirely finite summation ranges, define
$\beta_1=-\gamma$ and $\beta_j=(-1)^j(j-1)!\zeta(j)$ for $j\ge2$.
Set
\[
 \alpha_{j,-j}=(-1)^j(j-1)!,\quad \alpha_{1,0}=-\gamma,
 \quad
 \alpha_{j,k}=(-1)^{j+k}\frac{(j+k-1)!}{k!}\zeta(j+k)
 \quad(k\ge0,\ j+k\ge2).
\]
All other negative-index coefficients are zero. For a partition
$\lambda$ of $w$,
\begin{equation}\label{dq:finiteendpoint}
 \mathcal E(X_\lambda)=\prod_i\beta_{\lambda_i}
 -\sum_{\substack{
 k_i\in\{-\lambda_i,0,1,\ldots,w-\lambda_i\}\\
 k_1+\cdots+k_\ell=0}}
       \prod_i\alpha_{\lambda_i,k_i}.
\end{equation}
This supplies every endpoint constant at arbitrary degree without
leaving an unspecified finite-part operation.

The normal-form theorem therefore becomes the exact analytic identity
\begin{equation}\label{dq:integratednormalform}
 \operatorname{FP}\int_0^1F(\psi,\psi',\ldots)\,dx
 =\sum_{\lambda\in\mathscr B_N}
   [X_\lambda]\mathcal N(F)\,
   \operatorname{FP}\int_0^1\prod_i\psi^{(\lambda_i-1)}(x)\,dx
   +\mathcal E(\mathcal T(F)).
\end{equation}
Its proof uses finite algebra and \eqref{dq:endpointfunctional}; it
does not interchange divergent integrals of individual terms.

\subsection{Inverse-Gamma polynomials and the parity reduction}

Define formal differential polynomials by
\begin{equation}\label{dq:Arecurrence}
 A_0=1,\qquad A_{m+1}=X_1A_m-DA_m.
\end{equation}
On substituting $X_j=\psi^{(j-1)}$, these become
\[
 A_m(x)=(-1)^m\Gamma(x)\frac{d^m}{dx^m}\frac1{\Gamma(x)},
 \qquad
 \frac{\Gamma(x)}{\Gamma(x-t)}=
                 \sum_{m\ge0}A_m(x)\frac{t^m}{m!}.
\]
Write $N_m=\mathcal N(A_m)$ and $T_m=\mathcal T(A_m)$, with
$N_0=1,T_0=0,N_1=X_1,T_1=0$.

\begin{theorem}[Factor parity and an explicit all-degree recurrence]
\label{dq:parity}
Every monomial of $N_m$ has factor degree congruent to $m$ modulo two.
Thus its number of possible canonical coordinates is at most
\begin{equation}\label{dq:paritycount}
 r_m=\sum_{\ell\equiv m\ (2)}
             \bigl(p_\ell(m)-p_\ell(m-1)\bigr)\qquad(m\ge2).
\end{equation}
For $m\ge1$, the complete pair $(N_{m+1},T_{m+1})$ is computed
from $(N_m,T_m)$ by setting
\[
 (R,S)=(\mathcal N,\mathcal T)(X_1N_m-X_2T_m)
\]
and then
\begin{equation}\label{dq:pairrecurrence}
 N_{m+1}=R,\qquad
 T_{m+1}=S+X_1T_m-N_m-DT_m.
\end{equation}
All coefficient operations are finite and rational.
\end{theorem}

\begin{proof}
Here is a polynomial form of the adjoint argument, retaining its
boundary explicitly. Introduce a function $F$ with $F'=X_1$, and
write $p=e^F,q=e^{-F}$. The products below are differential
polynomials because the exponential factors cancel. The Lagrange
concomitant is
\[
 L_m=\sum_{j=0}^{m-1}(-1)^j p^{(j)}q^{(m-1-j)}.
\]
Termwise differentiation telescopes to
\[
 DL_m=pq^{(m)}+(-1)^{m-1}p^{(m)}q
       =(-1)^m A_m(X)-A_m(-X).
\]
Equivalently, put $B_j(X)=(-1)^jA_j(-X)$ and write
$L_m=\sum_{j=0}^{m-1}(-1)^jB_j(X)B_{m-1-j}(-X)$;
then the argument is an identity inside $\mathscr R$ and does not
require choosing an antiderivative $F$.

The normal-form projection respects factor degree, and replacing every
$X_j$ by $-X_j$ multiplies a degree-$\ell$ monomial by $(-1)^\ell$.
Applying the unique projection to the concomitant identity yields
$N_m(-X)=(-1)^mN_m(X)$, proving the parity assertion and
\eqref{dq:paritycount}. Finally substitute $A_m=N_m+DT_m$ into
\eqref{dq:Arecurrence} and use
$X_1DT_m=D(X_1T_m)-X_2T_m$. The result is
\[
 A_{m+1}=X_1N_m-X_2T_m+D(X_1T_m-N_m-DT_m),
\]
which proves \eqref{dq:pairrecurrence}.
\end{proof}

The values of $r_m$ for $m=2,\ldots,12$ are
\[
 1,1,2,2,3,3,5,6,8,9,13.
\]
These are formal upper bounds for the number of coordinates in
the particular Gamma normal form. No nonvanishing assertion for all
allowed coefficients is required. Some low-degree remainders are
\begin{align}
 N_4&=X_1^4-X_2^2,\notag\\
 N_5&=X_1^5-5X_1X_2^2,\notag\\
 N_6&=X_1^6-15X_1^2X_2^2+X_3^2,\label{dq:N4567}\\
 N_7&=X_1^7-35X_1^3X_2^2+7X_1X_3^2,\notag\\
 N_8&=X_1^8-70X_1^4X_2^2+28X_1^2X_3^2-7X_2^4-X_4^2.
 \notag
\end{align}
In particular the weight-six basis element $X_2^3$ is excluded by
parity. Its absence is a structural consequence, not a fitted numerical
cancellation.

\subsection{Appell transport of the reduced differential polynomials}

Let
\[
 P_j(Y)=j![t^j]\Gamma(1-t)e^{tY},\qquad
 C_m(x)=m![t^m]\frac{\Gamma(1-t)\Gamma(x)}{\Gamma(x-t)}.
\]
These are the established scalar and differential Gamma normalizers
from the incoming fifth-moment report. The operator $\partial_{X_1}$
commutes with $D$, and the repeated-largest normal-form space is
stable under $\partial_{X_1}$. Uniqueness gives
\[
 \partial_{X_1}N_m=mN_{m-1}.
\]
Consequently the canonical remainder of $C_m$ is obtained by expanding
$N_m$ as a polynomial in $X_1$ and replacing each power $X_1^j$ by
$P_j(\psi)$. Explicitly, with $K_j=N_j|_{X_1=0}$,
\begin{equation}\label{dq:Appelltransport}
 \mathcal N(C_m)=\sum_{j=0}^m\binom mj P_j(\psi)K_{m-j}.
\end{equation}
For example,
\begin{align}
 \mathcal N(C_6)&=P_6(\psi)-15P_2(\psi)(\psi')^2+(\psi'')^2,
 \label{dq:C6normal}\\
 \mathcal N(C_7)&=P_7(\psi)-35P_3(\psi)(\psi')^2
                                  +7P_1(\psi)(\psi'')^2.
 \label{dq:C7normal}
\end{align}
The primitive can be obtained by the same finite Gamma coefficient
convolution applied to the $T_j$; constant ambiguities do not affect
its endpoint functional.

Moreover, if $C_m=\mathcal N(C_m)+D\widetilde T_m$, then
\begin{equation}\label{dq:Appellendpointalgebra}
 \mathcal E(\widetilde T_m)\in
            \mathbb Q[\zeta(2),\ldots,\zeta(m-1)].
\end{equation}
It is homogeneous of weight $m-1$. To see this, put $y=\psi+\gamma$
and use
\[
 \Gamma(1-t)e^{-\gamma t}
  =\exp\left(\sum_{j\ge2}\frac{\zeta(j)}j t^j\right).
\]
The primitive is a differential polynomial in $y$ with these zeta
coefficients. Both $y(1)=0$ and the constant Laurent coefficient of
$y$ at zero are zero, so neither endpoint calculation introduces
$\gamma$. The finite formula \eqref{dq:finiteendpoint} proves
the assertion at every degree.

\subsection{The inherited spectral bridge in the new normal form}

Use the conventions
\[
 \frac1{\Gamma(1+u)}=\sum_{j\ge0}c_ju^j,\qquad
 \lambda_d=[u]\zeta(1+u,\underbrace{1,\ldots,1}_{d-1}),
 \quad d\ge1.
\]
The coefficient $[u]$ is a Laurent coefficient, not a derivative at
a pole. In particular $\lambda_1=-\gamma_1$ under the convention
$\zeta(1+u)=u^{-1}+\gamma-\gamma_1u+\cdots$.
The established Gamma integral bridge is
\begin{equation}\label{dq:inheritedCbridge}
 \int_0^1C_m(x)\,dx=m!(c_m-\lambda_{m-1}),\qquad m\ge2.
\end{equation}
Its complete integrand is continuous at zero. Combining this formula
with \eqref{dq:Appelltransport} and the exact endpoint functional
gives an all-degree finite recipe which isolates the single highest
height-one coefficient and specifies every remaining mixed integral.

For comparison, the equivalent inverse-Gamma bridge is
\begin{equation}\label{dq:inheritedAbridge}
 \mathcal L_m:=\operatorname{FP}\int_0^1 A_m(x)\,dx
 =m![t^m]\frac{\sin(\pi t)}{\pi t}
  -m!\sum_{d=1}^{m-1}(-1)^{m-d-1}c_{m-d-1}\lambda_d.
\end{equation}
This follows directly by multiplying the generating family $C_m$
by $1/\Gamma(1-t)$, retaining
$\operatorname{FP}\int C_0=1$ and
$\operatorname{FP}\int C_1=\gamma$, and using
$1/(\Gamma(1-t)\Gamma(1+t))=\sin(\pi t)/(\pi t)$.
Thus the new normal form changes the finite coordinate list and
retains the already established spectral normalization.

\section{A direct formula for coincident polygamma products}
\label{dq:bilinearsection}

The manuscript already gives general translated polygamma products
in first Hurwitz derivatives, in particular its formula
\texttt{shj:polyclosurezeta}. The following coincident, original-coordinate
finite part has the corresponding ordinary-zeta form. A direct proof
is supplied because a collision limit cannot in general be moved through
a finite-part prescription without checking its local normalization.

Put $H_0=0$ and $H_d=\sum_{j=1}^d1/j$.
For nonnegative integers $r,s$ with $d=r+s\ge1$, define
$\sigma_r=(-1)^{r+1}$ and $\sigma_s=(-1)^{s+1}$.

\begin{theorem}[Coincident bilinear polygamma finite parts]
\label{dq:bilinear}
In the original coordinate $x$,
\begin{align}
 &\operatorname{FP}\int_0^1\psi^{(r)}(x)\psi^{(s)}(x)\,dx\notag\\
 &\quad=d!\left\{(\sigma_r+\sigma_s)\zeta'(d+1)
 +\left[\sigma_r(H_d-H_s)+\sigma_s(H_d-H_r)\right]
                                      \zeta(d+1)\right\}.
 \label{dq:bilinearformula}
\end{align}
In particular, for every $r\ge1$,
\begin{equation}\label{dq:squareformula}
 \operatorname{FP}\int_0^1(\psi^{(r)}(x))^2\,dx
 =2(-1)^{r+1}(2r)!
    \left\{\zeta'(2r+1)+(H_{2r}-H_r)\zeta(2r+1)\right\}.
\end{equation}
\end{theorem}

\begin{proof}
At the pole $x=-n$, write $v=x+n$ and
\[
 \psi(-n+v)=-v^{-1}+a_0(n)+\sum_{k\ge1}a_k(n)v^k,
\]
where $a_0(n)=H_n-\gamma$ and
$a_k(n)=H_n^{(k+1)}+(-1)^{k+1}\zeta(k+1)$ for $k\ge1$.
In every case $\Delta a_k(n)=n^{-k-1}$.
Let $p_j(n)$ denote the coefficient of $v^{-j}$ in the product.
There is a constant highest coefficient at order $d+2$, and for
$1\le j\le d+1$ the other coefficients are
\begin{equation}\label{dq:bilinearpoles}
 p_j(n)=(d+1-j)!a_{d+1-j}(n)
 \left\{\mathbf1_{j\le r+1}\frac{\sigma_r r!}{(r+1-j)!}
       +\mathbf1_{j\le s+1}\frac{\sigma_s s!}{(s+1-j)!}\right\}.
\end{equation}
This is just multiplication of each leading pole by the other's
analytic Taylor series. In particular
\[
 \Delta p_1(n)=d!(\sigma_r+\sigma_s)n^{-d-1}.
\]

Here are the convergence and normalization details. The normalized
Mittag--Leffler expansion is
\[
 f(z)=P_0(z)+\kappa+
           \sum_{n\ge1}\{P_n(z)-P_n(0)\},\qquad
 P_n(z)=\sum_j\frac{p_j(n)}{(z+n)^j},
\]
where $\kappa$ is the local Laurent constant at zero. Its normal
convergence follows from $p_j(n)=O(1+\log n)$ and one subtraction.
The whole-circle contour argument eliminates the entire remainder:
integrate $z(f(w)-P_0(w))/(w(w-z))$ on $|w|=N+1/2$.
Digamma reflection and the differentiated cotangent reflection formulas
bound all factors by $O(1+\log N)$ on this circle, so the integral
tends to zero. Moreover
\[
 \kappa-\sum_{n=1}^NP_n(0)
       =\frac1{2\pi}\int_{-\pi}^{\pi}f((N+1/2)e^{i\theta})\,d\theta
       \longrightarrow0.
\]
For the last limit, outside an arc of length $O(\log N/N)$ at the
negative real axis, at least one differentiated digamma factor tends
to zero by the sectorial asymptotic and reflection. On the exceptional
arc the uniform logarithmic bound times its length tends to zero.
This uses $d\ge1$; the undifferentiated square is a separate case.

Integrate the normal expansion to one before passing to the cutoff.
Finite summation by parts in its higher-pole terms gives
\[
 \sum_{j\ge2}\sum_{n\ge1}
              \frac{\Delta p_j(n)}{(j-1)n^{j-1}},
\]
and the simple poles give $-\sum_{n\ge1}\Delta p_1(n)\log n$.
All terminal terms vanish: $p_j(N)/(N+1)^{j-1}\to0$ for $j\ge2$;
also $p_1(N)\log(N+1)\to0$, since it is zero when $d$ is odd,
and proportional to the decaying tail $a_d(N)$ when $d$ is even.
The grouped sums are absolutely convergent. They are precisely the
original-coordinate finite part, by integration before constant-term
extraction at zero.

The logarithmic sum is $d!(\sigma_r+\sigma_s)\zeta'(d+1)$.
For the higher poles use the finite identity
\begin{equation}\label{dq:binomialharmonic}
 r!\sum_{h=1}^r\frac{(d-h)!}{(r-h)!h}
                 =d!(H_d-H_{d-r}),\qquad0\le r\le d.
\end{equation}
After division by $d!$, this identity is obtained by integrating
\[
 \sum_{h=1}^r\binom rh t^{h-1}(1-t)^{d-h}
       =\frac{(1-t)^{d-r}-(1-t)^d}{t}
\]
from zero to one. The beta integrals on the left and the finite
geometric sum on the right give exactly \eqref{dq:binomialharmonic}.
Applying it once to each term in \eqref{dq:bilinearpoles} gives
the zeta coefficient in \eqref{dq:bilinearformula}. Setting $s=r$
proves \eqref{dq:squareformula}.
\end{proof}

The cases needed below are
\begin{equation}\label{dq:twosquares}
 \operatorname{FP}\int_0^1(\psi')^2\,dx=2\zeta(3)+4\zeta'(3),
 \qquad
 \operatorname{FP}\int_0^1(\psi'')^2\,dx=-28\zeta(5)-48\zeta'(5).
\end{equation}

\section{The canonical sixth and seventh digamma identities}
\label{dq:sixthsection}

Write $z_j=\zeta(j)$ in this section and put $y=\psi+\gamma$.
The sixth scalar polynomial is explicitly
\begin{align}
 P_6(\psi)={}&y^6+15z_2y^4+40z_3y^3
          +(45z_2^2+90z_4)y^2\notag\\
 &+(120z_2z_3+144z_5)y
     +15z_2^3+90z_2z_4+40z_3^2+120z_6.
 \label{dq:P6}
\end{align}
Define the one retained mixed moment
\begin{equation}\label{dq:M6}
 \widetilde M_6=\operatorname{FP}\int_0^1
                          (\psi(x)+\gamma)^2(\psi'(x))^2\,dx.
\end{equation}
This is a fully specified original-coordinate finite part, with an
ordinary convergent subtracted representation obtained by its finite
Laurent principal part.
More explicitly,
\begin{equation}\label{dq:M6principal}
 y^2(\psi')^2=x^{-6}-2z_3x^{-3}-z_4x^{-2}
                 +6(z_2z_3-z_5)x^{-1}+O(1).
\end{equation}
Consequently
\begin{align}
 \widetilde M_6={}&-\frac15+z_3+z_4\notag\\
 &+\int_0^1\left\{y^2(\psi')^2-x^{-6}+2z_3x^{-3}
         +z_4x^{-2}-6(z_2z_3-z_5)x^{-1}\right\}\,dx.
 \label{dq:M6ordinary}
\end{align}
The complete integrand in braces has a removable singularity at zero.

\begin{theorem}[A canonical sixth-degree bridge with one mixed moment]
\label{dq:canonicalsixth}
With the conventions already fixed,
\begin{align}
 \boxed{\begin{aligned}
 \operatorname{FP}\int_0^1P_6(\psi(x))\,dx
 ={}&720(c_6-\lambda_5)+15\widetilde M_6\\
 &+60z_2\zeta'(3)+48\zeta'(5)+90z_2z_3+270z_5.
 \end{aligned}}
 \label{dq:canonicalsixthformula}
\end{align}
Thus this scalar Gamma--Appell combination has only the highest
height-one coefficient $\lambda_5$ and one specified mixed moment,
apart from ordinary constants and ordinary zeta derivatives. This is
not an elimination or an independence theorem for $\widetilde M_6$.
\end{theorem}

\begin{proof}
For the inverse-Gamma polynomial, the exact primitive from the normal
algorithm is
\begin{align}
 T_6={}&-3\psi^5+20\psi^3\psi'-15\psi^2\psi''
       -15\psi(\psi')^2+6\psi\psi'''+9\psi'\psi''-\psi''''.
 \label{dq:T6}
\end{align}
Direct differentiation, or the finite recurrence, gives
\[
 A_6=\psi^6-15\psi^2(\psi')^2+(\psi'')^2+T_6'.
\]
For $C_6$, form the polynomial
\begin{align}
 \widetilde T_6={}&T_6[y]
  +15z_2(-2y^3+4y\psi'-\psi'')\notag\\
  &+40z_3(-\tfrac32y^2+\psi')-(45z_2^2+90z_4)y,
 \label{dq:TC6}
\end{align}
where $T_6[y]$ means \eqref{dq:T6} with every undifferentiated
$\psi$ replaced by $y$ and its derivatives unchanged. Then
\begin{equation}\label{dq:C6exact}
 C_6=P_6(\psi)-15(y^2+z_2)(\psi')^2+(\psi'')^2
                              +\widetilde T_6'.
\end{equation}
The two endpoint values are
\begin{align}
 \widetilde T_6(1)&=52z_2z_3+24z_5,\notag\\
 \operatorname{CT}_{x=0}\widetilde T_6(x)
                  &=112z_2z_3+266z_5,\notag\\
 \mathcal E(\widetilde T_6)&=-60z_2z_3-242z_5.
 \label{dq:C6endpoints}
\end{align}
They follow by substituting
$y=-x^{-1}+z_2x-z_3x^2+z_4x^3-z_5x^4+O(x^5)$,
or directly from \eqref{dq:finiteendpoint}. Apply finite parts to
\eqref{dq:C6exact}, use \eqref{dq:inheritedCbridge}, and then use
both evaluations in \eqref{dq:twosquares}. Every term in
\eqref{dq:canonicalsixthformula} follows, including the endpoint
contribution and the $z_2\zeta'(3)$ correction.
\end{proof}

For a fully expanded ordinary coefficient,
\begin{align}
 720c_6={}&\gamma^6-15\gamma^4z_2+40\gamma^3z_3
       +45\gamma^2z_2^2-90\gamma^2z_4\notag\\
 &-120\gamma z_2z_3+144\gamma z_5
       -15z_2^3+90z_2z_4+40z_3^2-120z_6.
 \label{dq:c6}
\end{align}

\subsection{The raw sixth power}

Let $Q_6=\operatorname{FP}\int_0^1\psi^6$ and
$M_6=\operatorname{FP}\int_0^1\psi^2(\psi')^2$.
The endpoint difference of \eqref{dq:T6} is
\begin{equation}\label{dq:rawsixthendpoint}
 \mathcal E(T_6)=60\gamma^3z_2-150\gamma^2z_3-105\gamma z_4
                              +210z_2z_3-242z_5.
\end{equation}
Combining this with \eqref{dq:inheritedAbridge} gives the complete
raw-power formula
\begin{align}
 Q_6={}&15M_6-720\lambda_5+720\gamma\lambda_4
          -360(\gamma^2-z_2)\lambda_3\notag\\
 &+120(\gamma^3-3\gamma z_2+2z_3)\lambda_2\notag\\
 &-30(\gamma^4-6\gamma^2z_2+8\gamma z_3
                       +3z_2^2-6z_4)\lambda_1\notag\\
 &-135z_6+48\zeta'(5)-60\gamma^3z_2+150\gamma^2z_3
             +105\gamma z_4-210z_2z_3+270z_5.
 \label{dq:rawsixth}
\end{align}
The term $-135z_6=-\pi^6/7$ is the entire contour constant of
the sixth inverse-Gamma bridge. It is not an endpoint correction.

\subsection{The next odd degree}

The algorithm also gives a compact seventh-degree extension. Define
\[
 W_7=\operatorname{FP}\int_0^1P_3(\psi)(\psi')^2\,dx,
 \qquad
 K_7=\operatorname{FP}\int_0^1P_1(\psi)(\psi'')^2\,dx.
\]
Then
\begin{equation}\label{dq:canonicalseventh}
 \boxed{\displaystyle
 \operatorname{FP}\int_0^1P_7(\psi)\,dx
 =5040(c_7-\lambda_6)+35W_7-7K_7+231z_3^2+2380z_6.}
\end{equation}
Here $P_1(\psi)=y$ and $P_3(\psi)=y^3+3z_2y+2z_3$.
For reproducibility, a primitive for $A_7-N_7$ is
\begin{align}
 T_7={}&-\tfrac72\psi^6+35\psi^4\psi'-35\psi^3\psi''
        -\tfrac{105}2\psi^2(\psi')^2+21\psi^2\psi'''\notag\\
 &+63\psi\psi'\psi''-7\psi\psi''''+14(\psi')^3
        -14\psi'\psi'''-\tfrac{21}2(\psi'')^2+\psi'''''.
 \label{dq:T7}
\end{align}
In addition to $T_6$ already displayed, the lower primitives needed
for the Gamma convolution are
\begin{align*}
 T_5&=-\tfrac52\psi^4+10\psi^2\psi'-5\psi\psi''
                         -\tfrac52(\psi')^2+\psi''',\\
 T_4&=-2\psi^3+4\psi\psi'-\psi'',\\
 T_3&=-\tfrac32\psi^2+\psi',\qquad T_2=-\psi.
\end{align*}
The primitive for $C_7-\mathcal N(C_7)$ is
\begin{align*}
 \widetilde T_7={}&T_7[y]+21z_2T_5[y]+70z_3T_4[y]\\
  &+(105z_2^2+210z_4)T_3[y]+(420z_2z_3+504z_5)T_2[y].
\end{align*}
The finite endpoint algorithm gives
\[
 \mathcal E(\widetilde T_7)
 =-448z_2z_4-231z_3^2-1596z_6
 =-231z_3^2-2380z_6.
\]
Combining it with \eqref{dq:C7normal} and the inherited Gamma
bridge proves \eqref{dq:canonicalseventh}. The certificate contains
the full polynomial identity, so this seventh-degree formula does
not rely on a numerical extrapolation from degree six.

\subsection{Proof status, verification, and remaining question}

The accompanying exact program checks every partition at every weight
through twelve: it verifies the decomposition identity, the vanishing
normal form of every total derivative, the bidegree count, the adjoint
concomitant identity, and the parity restriction. It also independently
expands the local Laurent polynomials for the displayed sixth and
seventh endpoint constants. These finite checks corroborate the
proofs; the terminating algorithm and the dimension theorem themselves
hold at all degrees by the arguments above.

This resolves the formal integration-by-parts classification requested
in the incoming programme and gives complete degree-six and degree-seven
applications. The next separate task is to determine the image of this
formal quotient in an explicitly specified family of harmonic spectral
coordinates. Its kernel may contain additional special-function
relations. Neither the partition count nor the parity count can be
used to assert arithmetic independence of the remaining constants.

% Suggested bibliography entry for the surrounding article:
% \bibitem{OlverShakiban1978}
% P. J. Olver and C. Shakiban, A resolution of the Euler operator. I,
% Proc. Amer. Math. Soc. 69 (1978), 223--229.
% DOI: 10.1090/S0002-9939-1978-0486822-8.
% Primary author-hosted PDF: https://www-users.cse.umn.edu/~olver/s_/reo.pdf
